
\documentclass[../../main.tex]{subfiles}
\begin{document}

% 年度の大きな表示
\begin{center}
  {\fontsize{48pt}{58pt}\selectfont \textbf{1984}\normalsize\textbf{年度}}
\end{center}

\vspace{10mm}

% 本文

\vspace{15mm}

% 出題テーマと難易度の表
\noindent\textbf{出題テーマと難易度}

\vspace{5mm}

\begin{center}
  \shadowbox{%
    \begin{tabular}{|c|p{8cm}|c|}
      \hline
      \textbf{問題番号} & \textbf{テーマ} & \textbf{難易度} \\
      \hline
      第1問 & 整数問題 & \\
      \hline
      第2問 & 整数解の個数 & \\
      \hline
      第3問 & 双曲線と面積 & \\
      \hline
      第4問 & 定積分の最小値 & \\
      \hline
      第5問 & 三角関数と面積 & \\
      \hline
    \end{tabular}%
  }
\end{center}

\vspace{15mm}

% 解答
\noindent\textbf{解答}

\vspace{5mm}

\begin{center}
  \shadowbox{%
    \renewcommand{\arraystretch}{1.6}
    \begin{tabular}{|c|c|p{8cm}|}
      \hline
      \textbf{問題番号} & \textbf{小問} & \textbf{答え} \\
      \hline
      \multirow{2}{*}{第1問} & (1) & 証明問題 \\
      \cline{2-3}
                             & (2) & $(a,b)=(2^n,2^n),\ (2\cdot3^n,3^n),\ (3^n,2\cdot3^n)$ \\
      \hline
      \multirow{2}{*}{第2問} & $n$偶数 & $\dfrac{1}{8}n^2+\dfrac{3}{4}n-2$個 \\
      \cline{2-3}
                             & $n$奇数 & $\dfrac{(n-1)(n+1)}{8}$個 \\
      \hline
      \multirow{2}{*}{第3問} & (1) & 証明問題 \\
      \cline{2-3}
                             & (2) & 証明問題 \\
      \hline
      第4問 & - & $a=\log\dfrac{e+1}{2}$ \\
      \hline
      第5問 & - & $T=-\dfrac{1}{2}\log\dfrac{\sqrt{5}-1}{2}-\dfrac{1}{8}(\sqrt{5}-1)^2$ \\
      \hline
    \end{tabular}%
    \renewcommand{\arraystretch}{1.0}
  }
\end{center}

\vspace{15mm}

% 解答時間
\noindent\textbf{試験形態}

\vspace{5mm}

\begin{center}
  \shadowbox{%
    \begin{tabular}{p{8cm}}
      （要確認）
    \end{tabular}%
  }
\end{center}

\end{document}
